\input{../tex-inputs/latex-leseansicht-vorspann}

\section[Arthur und Olga Schnitzler an Gustav Schwarzkopf, 18. 4. 1911]{L04446 Arthur und Olga Schnitzler an Gustav Schwarzkopf, 18. 4. 1911}
\nopagebreak\mylabel{L04446v}
\rehead{ }\normalsize\beginnumbering\briefempfaengerindex{Schwarzkopf, Gustav@\textsc{Schwarzkopf, Gustav}!zzzSchnitzler, Olga@\emph{von Olga Schnitzler}!1911-04-181@{18. 4. 1911}|(be}\briefempfaengerindex{Schwarzkopf, Gustav@\textsc{Schwarzkopf, Gustav}!zzzSchnitzler, Arthur@\emph{von Arthur Schnitzler}!1911-04-181@{18. 4. 1911}|(be}
\toendnotes[C]{\smallbreak\pagebreak[2]}
\correspDesc{Versand  durch Arthur Schnitzler, Olga Schnitzler am 18. 4. 1911 in Menton
\newline{}Übermittlung  am 18. 4. 1911 in Menton
\newline{}Erhalt  durch Gustav Schwarzkopf im Zeitraum [19. 4. 1911 – 23. 4. 1911?] in Wien}\toendnotes[C]{\smallbreak}
\Standort{DLA, A:Schnitzler, HS.1985.1.1897.}
\physDesc{Bildpostkarte, 235 Zeichen
\newline{}Handschrift Arthur Schnitzler: Bleistift, deutsche Kurrent
\newline{}Handschrift Olga Schnitzler: Bleistift, lateinische Kurrent
\newline{}Versand: Stempel: »\nobreak{}\oindex{Menton@\textbf{Menton}|pwk}Menton\hspace*{1em}Alpes Maritimes, 18–4 11, 2\textsuperscript{e} 15\nobreak{}«.  }\toendnotes[C]{\smallbreak}\pstart{}{\pb}\textsc{Autriche\oindex{Österreich@\textbf{Österreich}|pw}}\pend{}\pstart{}Herrn \textsc{Gustav Schwarzkopf}\pend{}\pstart{}Wien I\oindex{I., Innere Stadt@\textbf{I., Innere Stadt}|pw}\pend{}\pstart{}\textsc{Tiefer Graben 23}\oindex{Wien@\textbf{Wien}!I., Innere Stadt@\textbf{I., Innere Stadt}!Tiefer Graben 23@\textbf{Tiefer Graben 23}|pw}\pend{}{\bigskip}
\pstart
           \noindent{}{\pb}\textcolor{gray}{\textbf{MENTON\oindex{Menton@\textbf{Menton}|pw}. – \begin{otherlanguage}{french}LA VILLE ET LE QUAI\end{otherlanguage}}}\pend
           \vspace{1em}
\pstart
           {\pb}18. 4. 911{ }\textsc{Mentone\oindex{Menton@\textbf{Menton}|pw}}{\\}\textsc{Hotel National\oindex{Grand Hôtel National [Menton]@\textbf{Grand Hôtel National [Menton]}|pw}}. –\pend
           \vspace{0.5em}
\pstart
           Viele herzliche Grüße. Wir haben das{ }ſchönſte Wetter in 
      dieſer herrlichen
      Gegend. – \textsc{Liesl\pwindex{Steinrück, Elisabeth 19.\,11.\,1885 – 7.\,4.\,1920 Partenkirchen@\textsc{Steinrück, Elisabeth} (19.\,11.\,1885 – 7.\,4.\,1920 Partenkirchen)|pw}}
      gehts leidlich, wir 
      fahren über Par\textcolor{gray}{t}enk\oindex{Partenkirchen@\textbf{Partenkirchen}|pw}
      zurück. –\pend
           
\pstart
           Ihr \spacefill\mbox{A.}{\\[\baselineskip]}\spacefill\mbox{{[}hs. O. Schnitzler:{]} \label{T_L04446-1v}\edtext{OlgaS.}{\lemma{\textnormal{\emph{OlgaS.}}}\Cendnote{\textnormal{Olga Schnitzlers\pwindex{Schnitzler, Olga 17.\,1.\,1882 Wien – 13.\,1.\,1970 Lugano@\textsc{Schnitzler, Olga} (17.\,1.\,1882 Wien – 13.\,1.\,1970 Lugano)|pwk} Unterschrift befindet sich um 90° gedreht auf einer freien Fläche neben dem Text.}}}\label{T_L04446-1}}\pend
           \leftskip=0em{}\selectlanguage{ngerman}\endnumbering\briefempfaengerindex{Schwarzkopf, Gustav@\textsc{Schwarzkopf, Gustav}!zzzSchnitzler, Olga@\emph{von Olga Schnitzler}!1911-04-181@{18. 4. 1911}|)be}\briefempfaengerindex{Schwarzkopf, Gustav@\textsc{Schwarzkopf, Gustav}!zzzSchnitzler, Arthur@\emph{von Arthur Schnitzler}!1911-04-181@{18. 4. 1911}|)be}\mylabel{L04446h}
\begin{anhang}
\end{anhang}\newcommand{\dateiname}{L04446}\newcommand{\titel}{Arthur und Olga Schnitzler an Gustav Schwarzkopf, 18. 4. 1911}\newcommand{\editorInnen}{Herausgegeben von Jahnke, SelmaMüller, Martin Anton}\input{../tex-inputs/latex-leseansicht-abspann}
